\section{Conclusion}

We have presented Quantification, a fully automated pipeline for scratch-wound healing
quantification that combines variance-based single-frame detection with a biologically
motivated monotonic closure constraint.
The pipeline requires no training data, processes standard phase-contrast microscopy images in
linear time, and produces smooth, strictly non-increasing wound area trajectories.
Validation on 325 images from a pharmacological study in immortalized myoblasts demonstrated
98.9\% detection accuracy and a mean Dice coefficient of 0.91 against expert annotations.
Application to 282 validated trajectories revealed that ALK5 inhibition suppressed wound
closure kinetics while candesartan had no appreciable effect, consistent with a dominant role
for canonical TGF-$\beta$ signaling in myoblast migration.
A three-tier quality-control system ensures robust performance, and parallel trajectory
processing enables scalable analysis of high-content experiments.
The mathematical formulation is presented in full, all parameters are documented, and the
implementation is released as open-source Python code.
We anticipate that Quantification will reduce the time and subjectivity currently associated
with scratch-wound quantification, enabling researchers to focus on biological interpretation
rather than manual measurement.

\bigskip
\textbf{Data and Code Availability.}
The Quantification segmentation pipeline is implemented in Python~3.11 and is available at
[repository URL].
Raw microscopy images and processed quantitative data are available at [data repository].
All analysis scripts, configuration files, and parameter documentation are included in the
repository.
